\chapter{Glossing Abbreviations and Notation}\label{app:abbreviations}

\section{Glossing abbreviations}
Interlinear examples distinguish lexical meanings from grammatical categories,
whose abbreviations are set in small capitals. A hyphen separates affixes from
their hosts, while a full stop separates grammatical values expressed by a
single form. Zero exponence is represented by $\varnothing$ where the analysis
requires an unpronounced exponent. Comparative examples retain the
case and tense labels of the language illustrated; these labels do not
identify additional Iridian inflectional categories
(§\,\ref{sec:lexical-slot-schemata}).

\begingroup
\small
\renewcommand{\arraystretch}{1.05}
\begin{longtable}{@{}lp{0.74\textwidth}@{}}
\toprule
Abbreviation & Meaning \\
\midrule
\endfirsthead
\toprule
Abbreviation & Meaning \\
\midrule
\endhead
\bottomrule
\endfoot
\mk{1, 2, 3} & first, second, third person \\
\mk{adn} & adnominal form \\
\mk{all} & allative \\
\mk{assert} & assertive particle \\
\mk{assoc} & associative plural \\
\mk{attr} & attributive \\
\mk{cl} & classifier \\
\mk{comp} & comparative \\
\mk{decl} & declarative \\
\mk{exh} & exhortative \\
\mk{inch} & inchoative \\
\mk{prs} & present \\
\mk{pst} & past \\
\mk{state} & state-forming derivation \\
\mk{acc} & accusative \\
\mk{dat} & dative \\
\mk{gen} & genitive \\
\mk{ins} & instrumental \\
\mk{loc} & locative \\
\mk{nom} & nominative \\
\mk{obl} & oblique \\
\mk{q} & interrogative \\
\mk{dir} & direct case \\
\mk{trs} & transitive case \\
\mk{ind} & indirect case \\
\mk{voc} & vocative \\
\mk{pf} & perfective \\
\mk{ret} & retrospective \\
\mk{cont} & continuous \\
\mk{prog} & progressive \\
\mk{ctpl} & contemplative \\
\mk{antip} & antipassive \\
\mk{appl} & applicative \\
\mk{caus} & causative \\
\mk{imp} & imperative \\
\mk{proh} & prohibitive \\
\mk{hort} & hortative \\
\mk{irr} & irrealis \\
\mk{n.irr} & negative irrealis \\
\mk{seq} & sequential converb \\
\mk{circ} & circumstantial converb \\
\mk{sim} & simultaneous converb \\
\mk{purp} & purposive converb \\
\mk{cncs} & concessive converb \\
\mk{cond} & open conditional converb \\
\mk{e.cond} & expected conditional converb \\
\mk{cv} & converb (generic) \\
\mk{disc} & discourse connective \\
\mk{neg} & negation \\
\mk{def} & definiteness marker \\
\mk{s.anch} & speaker-anchored (Anch\textsuperscript{0}[+anch,\,\textminus{}shr]) \\
\mk{indef} & indefinite modifier \\
\mk{pl} & plural \\
\mk{sg} & singular \\
\mk{dem} & demonstrative \\
\mk{prox} & proximal \\
\mk{med} & medial \\
\mk{dist} & distal \\
\mk{anim} & animate \\
\mk{inan} & inanimate \\
\mk{refl} & reflexive \\
\mk{quot} & quotative \\
\mk{compl} & complementiser \\
\mk{lnk} & linker \\
\mk{exst} & existential \\
\mk{cop} & copula \\
\mk{neg.cop} & negative copula \\
\mk{neg.exst} & negative existential \\
\mk{nmlz} & nominaliser \\
\end{longtable}
\endgroup

\section{Forms and analytical notation}
Iridian forms are italicised, translations appear in single quotation marks,
and phonetic transcriptions are enclosed in square brackets; angle brackets
identify orthographic representations. An asterisk marks an ungrammatical form
or, in an explicitly historical discussion, a reconstructed form.
Functional heads and projections retain their symbolic names: \synhead{DP}
for the nominal domain, \synhead{AspP} for aspect, \synhead{VoiceP} for
participant structure, and \synhead{PolP} for clause typing. Their structural
relations are set out in §\,\ref{sec:grammatical-architecture}, while nominal
inflection is treated in Chapter~\ref{ch:nominal-morphology} and anchoring
in Chapter~\ref{ch:nominal-syntax}.
